\section{The record in \LaTeX, and three independent checkers}\label{sec:latex}

Mathematical texts are written in \LaTeX, and language models write \LaTeX{} readily. The natural place to keep a record is therefore the
document itself: each statement is declared where it appears in the text, and the \TeX{} compiler applies the rules while it typesets.
The package \texttt{qam.sty} does this. It uses only the \LaTeX{} kernel (expl3, \LaTeX{} 2022-06 or later) and depends on no
application.

\subsection{Declaring the example}

The following block is part of the source of this article. It records the interval prediction of Section~\ref{sec:example}. The frames
and the map were declared in Section~\ref{sec:frames}; the optional first argument of each command is the symbol under which the entry
is displayed.

\begin{lstlisting}
\qamthesis[\mathcal P]{preparation}{Lab}{conditional}
  {E=2e_x, B=t e_y+u e_z, t^2+u^2=1, c=1, boost v=1/4 along z}
\qaminput[t_I]{measurement}{Lab}{conditional}
\qamconvert[p]{range-rule}{conditional}{lorentz-z}{preparation}
\qamderive[\mathcal P_I]{measured-preparation}{conditional}{preparation,measurement}
\qamconvert[p_I]{narrowed-range}{conditional}{lorentz-z}{measured-preparation}
\end{lstlisting}
\qamthesis[\mathcal P]{preparation}{Lab}{conditional}
  {E=2e_x, B=t e_y+u e_z, t^2+u^2=1, c=1, boost v=1/4 along z}
\qaminput[t_I]{measurement}{Lab}{conditional}
\qamconvert[p]{range-rule}{conditional}{lorentz-z}{preparation}
\qamderive[\mathcal P_I]{measured-preparation}{conditional}{preparation,measurement}
\qamconvert[p_I]{narrowed-range}{conditional}{lorentz-z}{measured-preparation}


\begin{figure}[t]
\centering
\begin{tikzpicture}[node distance=10mm and 14mm,
  box/.style={draw,rounded corners=2pt,align=center,font=\small,minimum width=30mm,minimum height=9mm,fill=white},
  inp/.style={box,dashed},
  arr/.style={-{Stealth[length=2mm]},thick}]
\node[box] (P) at (0,0) {$\mathcal P$: preparation\\[-1pt]\footnotesize thesis, conditional};
\node[inp] (T) at (6.4,0) {$t_I$: measurement\\[-1pt]\footnotesize input, conditional};
\node[box] (PI) at (4.2,-1.9) {$\mathcal P_I$: restricted preparation\\[-1pt]\footnotesize derive, conditional, inputs $\{t_I\}$};
\node[box] (p) at (0,-4.6) {$p$: general rule\\[-1pt]\footnotesize convert, conditional};
\node[box] (pI) at (4.2,-4.6) {$p_I$: narrowed prediction\\[-1pt]\footnotesize convert, conditional, inputs $\{t_I\}$};
\draw[arr] (P) -- (PI); \draw[arr] (T) -- (PI);
\draw[arr] (P) -- node[left,font=\footnotesize]{$\Lambda_z$} (p);
\draw[arr] (PI) -- node[right,font=\footnotesize]{$\Lambda_z$} (pI);
\begin{scope}[on background layer]
\node[fill=black!5,rounded corners,fit=(P)(T)(PI),inner sep=4mm,label={[font=\footnotesize]north west:frame Lab}] {};
\node[fill=black!10,rounded corners,fit=(p)(pI),inner sep=4mm,label={[font=\footnotesize]south west:frame Boosted}] {};
\end{scope}
\end{tikzpicture}
\caption{The record of the interval prediction. Arrows point from a premise to what is obtained from it; the two arrows labelled
$\Lambda_z$ are conversions through the declared boost map, the only way from the Lab frame to the Boosted frame. The dashed box is an
input. Its key travels down to every entry below it, so $p_I$ carries it, while $p$ does not.}
\label{fig:record}
\end{figure}

\noindent Figure~\ref{fig:record} shows the resulting graph. The thesis $\mathcal P$ is the preparation~\eqref{eq:prep} with $v=\frac14$; as a thesis, its wording is part of its
hash, so a change of the set-up changes every hash below it; the input $t_I$ is the measured interval
$t\in[-\frac14,\frac14]$; $p$ is the general rule~\eqref{eq:rho-boosted} in the Boosted frame; $\mathcal P_I$ is the preparation
restricted by the measurement; and $p_I$ is the narrowed prediction~\eqref{eq:interval}. The author states the status of each entry.
The frame and the inputs of a derivation are not stated but computed from the premises, so they cannot be stated wrongly.

From these declarations the package typesets a compact notation. \verb|\qam{narrowed-range}| and \verb|\qam{range-rule}| give
\begin{center}
\qam{narrowed-range}\qquad\qquad \qam{range-rule}
\end{center}
Here $\Lambda_z$ is the symbol declared for the boost map, the left superscript is the frame, ``$\Leftarrow$'' reads ``derived
from'', $R(w)$ is the readout of a witness $w$ (a value put in by hand), and $W$ lists the witnesses an entry rests on. Similarly,
\verb|\qamshow{narrowed-range}| gives the short form \qamshow{narrowed-range}. One reads off that $p_I$ was carried from Lab to Boosted,
rests on the measured input $t_I$ and holds at most conditionally, while the general rule $p$ rests on no input.

The package also prints a ledger of all entries. For the five entries above it reads:
\begin{center}
\small
\begin{tabular}{@{}lll@{}}
\toprule
key & QAM notation & hash\\
\midrule
\texttt{preparation} & \qam*{preparation} & \texttt{\qamhash{preparation}}\\
\texttt{measurement} & \qam*{measurement} & \texttt{\qamhash{measurement}}\\
\texttt{range-rule} & \qam*{range-rule} & \texttt{\qamhash{range-rule}}\\
\texttt{measured-preparation} & \qam*{measured-preparation} & \texttt{\qamhash{measured-preparation}}\\
\texttt{narrowed-range} & \qam*{narrowed-range} & \texttt{\qamhash{narrowed-range}}\\
\bottomrule
\end{tabular}
\end{center}
The full record of the article, including the entries of Section~\ref{sec:frames}, is reproduced in Appendix~\ref{app:record}.

\subsection{What happens when a rule is broken}

By default a violation stops the compiler. This is deliberate: a document containing a promotion should not come out as a PDF. Had the
last line above declared the prediction \st{proved}, typesetting would have stopped with
\begin{lstlisting}
! Package qam Error: Promotion: `narrowed-range' is declared `proved', but its
(qam)                weakest premise `measured-preparation' is only
(qam)                `conditional'.
\end{lstlisting}
A declaration that the prediction is independent of inputs, \verb|\qamindependent{narrowed-range}|, stops with
\begin{lstlisting}
! Package qam Error: `narrowed-range' is cited as independent of inputs, but
(qam)                it rests on: measurement.
\end{lstlisting}
and a direct comparison of the prediction with the preparation,
\begin{center}\verb|\qamcompare{narrowed-range}{preparation}|,\end{center}
stops with
\begin{lstlisting}
! Package qam Error: Frame mixing: `narrowed-range' (frame Boosted) and
(qam)                `preparation' (frame Lab) are compared without a
(qam)                declared map.
\end{lstlisting}
These messages were produced by compiling the three variants; they are quoted, not paraphrased. Further messages cover keys, frames and
maps that are used before they are declared or declared twice, unknown status names, a conversion whose map starts in the wrong frame,
and a derivation without premises. A warning mode (\texttt{mode=warning}) turns errors into warnings so that the package can be
introduced into an existing text; the record exported in that mode is then rejected by the Lean checker, so the violation cannot pass
unnoticed.

\subsection{The exported record}

With the option \texttt{export}, the package writes every entry as one line of a plain text file. For the example, the lines of the
prediction and its ancestors read (hashes shortened, lines wrapped)
\begin{lstlisting}
map lorentz-z : Lab -> Boosted
preparation := thesis; frame Lab; status conditional; inputs {};
    hash 839F0699...; sources {}; statement E=2e_x, B=t e_y+u e_z, ...
measurement := <6,measurement>; frame Lab; status conditional;
    inputs {measurement}; hash C2DEF159...; sources {}
measured-preparation := derive(preparation,measurement); frame Lab;
    status conditional; inputs {measurement}; hash 16E7B1C0...;
    sources {839F0699...,C2DEF159...}
narrowed-range := convert(lorentz-z,measured-preparation); frame Boosted;
    status conditional; inputs {measurement}; hash 3F78DEBD...;
    sources {16E7B1C0...}
\end{lstlisting}
An input is written as \texttt{<6,measurement>}: the format tags every value put in by hand as a \emph{witness}, one of a small
set of tagged object kinds that the request language of Section~\ref{sec:requests} also uses. Appendix~\ref{app:record} reproduces the full exported record of this article. The exported file is the interface to every other
checker. Checking does not depend on trusting the \TeX{} package.

\subsection{Three checkers}

Three programs, written independently of one another, check an exported record. Two of them apply the rules of
Definition~\ref{def:rules}; the third checks the hashes:
\begin{enumerate}[leftmargin=1.5em]
\item the \TeX{} package, while the document is typeset;
\item \texttt{qamcheck}, a program written in Lean~4, on the exported file. It performs the one-pass check of
Section~\ref{sec:properties}, including the hash links, and its correctness properties are proved in Lean (Appendix~\ref{app:lean});
\item \texttt{qamhash.py}, a short Python script that recomputes every entry hash, the record hash and the MD5 of the source file.
\end{enumerate}
Lean does not compute MD5; the script does. Lean checks that the links between the recorded hashes are consistent. For a release, a
manifest of SHA-256 sums over the source, the PDF and the record is written in addition.

A checker that only ever sees good records is not tested. The repository therefore contains deliberately broken records that must be
rejected: a promotion, a frame mixing, a false declaration of independence, a conversion whose map starts in the wrong frame, a broken hash
link, an entry whose statement was edited after export, and an input that was deleted by hand. The last cannot be produced by the \TeX{}
package at all. It ensures that the Lean checker does more than repeat what the package already prevents. The \TeX{} package has its
own test suite of 34 small documents, each of which must either compile or stop with a specific error.
